\section{Introduction}\label{sec:introduction}

For a fixed field $\F$, the matrix multiplication exponent
$\omega=\omega_\F$ is the infimum of the real numbers $\tau$ for which
two $n\times n$ matrices over $\F$ can be multiplied using
$O(n^{\tau+o(1)})$ scalar additions and multiplications.  This exponent
measures the asymptotic arithmetic cost of a central bilinear operation
and of many algorithms built from it.  The field is fixed throughout;
constants in an arithmetic algorithm may depend on the field and on its
chosen finite construction.

Strassen's seven-product algorithm for $2\times2$ matrices gave the
subcubic exponent $\log_2 7$~\cite{Strassen1969}.  Subsequent advances
used tensors to record bilinear operations: tensor rank counts the
products of linear forms needed to realize an operation.  Approximate
bilinear identities and their conversion to exact algorithms
\cite{Bini1980}, Sch\"onhage's asymptotic sum inequality
\cite{Schonhage1981}, and Strassen's laser method
\cite{Strassen1987} established the main ingredients of this approach.
The sum inequality converts a low-rank direct sum of matrix tensors into
an exponent bound.  The laser method constructs such direct sums by
restricting powers of a tensor with low asymptotic rank.  Coppersmith and
Winograd combined these ideas with progression-free sets and introduced
the tensor family used here~\cite{CoppersmithWinograd1990}.

Passing to a larger tensor power gives more freedom in choosing the
intermediate pieces before extracting matrix tensors from them.
Stothers analyzed the fourth power, in work subsequently published with
Davie~\cite{Stothers2010,DavieStothers2013}.  Vassilevska Williams
analyzed the eighth power, and Le Gall continued through the
thirty-second power~\cite{VassilevskaWilliams2012,LeGall2014}.
Alman and Vassilevska Williams then reduced the extraction loss that
arises when prescribed marginal distributions permit several joint
distributions, obtaining $\omega<2.3728596$
\cite{AlmanWilliams2024}.

A further source of loss is requiring entire variable blocks to be
disjoint when the final matrix tensors need only disjoint subsets of
those blocks.  Duan, Wu, and Zhou used asymmetric hashing to reduce this
\emph{combination loss} in the recursive analysis
\cite{DuanWuZhou2023}.  Vassilevska Williams, Xu, Xu, and Zhou allowed
finer sharing of variables and extended recovery from deletions on all
three sides to the restricted tensor products needed in the recursion
\cite{VXXZ2023}.  Alman, Duan, Vassilevska Williams, Xu, Xu, and Zhou
then introduced successive compatibility tests that treat all three
sides differently~\cite{MoreAsymmetry2026}.  These tests underlie the
extraction studied here.

For a recent numerical comparison, Dupont et al.\ obtained
$\omega<2.371177$ by optimizing the eighth-power instance of this
asymmetric framework, combining modern optimization with AlphaEvolve
\cite{Dupont2026}.  Version~3 of \emph{More Asymmetry Yields Faster
Matrix Multiplication} records that improvement and credits the new
parameters to Dupont et al.\ \cite[Table~1]{MoreAsymmetry2026}.
We prove the following bound over every fixed field.

\begin{theorem}\label{thm:main}
Over every fixed field $\F$, the arithmetic exponent of square matrix
multiplication satisfies
\[
 \omega_\F<2.371054886006746<2.371054887.
\]
\end{theorem}

The construction gives exact arithmetic algorithms, including in positive
characteristic.  Its finite choices are independent of the input matrices.
The theorem concerns asymptotic arithmetic cost; the counting arguments
and parameter choices do not give bit-complexity or practical crossover
estimates.

\subsection{Proof overview}

We use the Coppersmith--Winograd tensor with parameter $5$, denoted
$\CW_5$.  Section~\ref{sec:foundations} defines it and proves that the
rank of its $m$th tensor power is at most $7^m$ times a polynomial in
$m$, over the original field.  A \emph{strand of length $\ell$} is a
copy of $\CW_5^{\otimes\ell}$.  On each side the variables of $\CW_5$
are indexed by $0,\ldots,6$, with weights $0,1,1,1,1,1,2$.
A strand variable is a word of these indices, and a \emph{shape} records
the weight sums on the three sides.  We start with strands of
length $8$ and split them successively into lengths $4$, $2$, and $1$.
At each depth we restrict by shape.  A shape with a zero coordinate
already gives a matrix multiplication tensor; positive shapes pass to
the next split.

\paragraph{Combining capacities across depths.}
An extraction restricts a tensor to a direct sum of specified tensor
products.  To make the outputs disjoint, its three side-assignment tests
each impose a bound on the guaranteed logarithmic number of outputs;
we call these bounds its \emph{directional capacities}.  The guarantee
uses their minimum.  The constituent-stage analysis in
\cite[Theorem~6.4]{MoreAsymmetry2026} already combines differently
prescribed populations at one depth before taking this minimum.
Our joint step also combines populations at different depths, provided
their priority orders agree on the physical tensor sides.  For their
weighted capacity vectors $v_r\in\mathbb R^3$, the benefit of pooling
across depths is expressed by
\[
 \sum_r\min_W v_{r,W}\ \leq\ \min_W\sum_r v_{r,W}.
\]
Thus a shortage on one side at one depth can be offset by capacity from
another depth.

Section~\ref{sec:joint} proves this joint extraction by following the
actual order of the tests.  On the first side, variables are grouped by
their sequences of strand or half-strand weight sums; these groups are
the \emph{coarse blocks}.  Each retained coarse block is assigned to a
unique output.  Compatibility with that assignment then determines which
complete variables on the second and third sides
can be retained.  The result is a direct sum of the intended products
with some variables deleted.  Populations retain their ancestry and
physical placement throughout, so each subsequent split receives its
own prescribed distributions.

\paragraph{Recovering outputs with inherited restrictions.}
Some variables are missing because an earlier extraction imposed type
conditions, meaning conditions on empirical distributions of variable
statistics.  Others are deleted to resolve competing assignments in the
current extraction.  Earlier recovery arguments already treat
three-side deletions and inherited type conditions
\cite[Section~7]{VXXZ2023}, \cite[Section~6.5]{MoreAsymmetry2026}.
Section~\ref{sec:recovery} proves the form needed for our joint step.
With $N$ the large strand-count parameter, the extraction selects outputs
in which every orbit of complete variables under the allowed position
permutations loses at most $C/N$ of its variables, uniformly over the orbits.  Here $C$ depends only on
the fixed population data and tolerances.  Independent symmetry shifts
then cover the ideal tensor.  Partitioning variables by their presence
in those shifts reconstructs every coefficient exactly once, over the
original field.  Only a subexponential number of independent input copies
is needed.  The recovered products can then be subdivided and used in
further extractions.

\paragraph{Staggering and the final rank comparison.}
A \emph{lot} contains $N$ length-$8$ strands.  We stagger finitely many
lots so that, at an interior time step, one lot is at each of the three
splitting stages.  Stage priorities balance the summed directional
capacities.  The work is partitioned among the six physical axis orders;
one joint extraction acts on each part, and the six output counts
multiply.  Section~\ref{sec:stagger} constructs this schedule, shown in
Figure~\ref{fig:pipeline}, and computes its finite boundary loss.
Section~\ref{sec:leaf-volumes} counts the common oriented dimensions
of the terminal matrix tensors and combines these counts with the
schedule to prove the total yield and rank comparison.

Write $H_0$ for the initial logarithmic yield, $C_*$ for the contribution
of the three staggered stages, and $S_*$ for the terminal logarithmic
matrix volume, all normalized per initial length-$8$ strand.  The
direct-sum rank inequality then gives
\begin{equation}\label{eq:intro-constraint}
 \omega_\F\leq\frac{3(8\log7-H_0-C_*)}{S_*}.
\end{equation}
For each fixed number of lots, we first take the divisible strand count
$N$ to infinity and remove approximation losses; the number of lots
tends to infinity last.  Thus the boundary cost vanishes without a
uniform assertion for a growing schedule.  The explicit rational data
in Section~\ref{sec:parameters} and Appendix~\ref{app:parameter-tables}
certify $S_*>0$ and place the ratio in
\eqref{eq:intro-constraint} below the bound in Theorem~\ref{thm:main}.
The complete parameter verifier and its rational-logarithm interval
certificate accompany the source.  The proof requires a single feasible
choice of parameters, with no assertion that this choice is optimal.
